\documentclass[11pt]{article}
\usepackage[utf8]{inputenc}
\usepackage[margin=2cm]{geometry}
\usepackage{hyperref}

\begin{document}
\noindent
November 13, 2026

\vspace{0.6em}
\noindent
To the Editor-in-Chief and Editorial Board of \textit{Renewable Energy}

\vspace{0.6em}
\noindent
Dear Editor,

\vspace{0.5em}
\noindent
We are pleased to submit our original research article entitled \textit{``Scale-Bridging Techno-Economic, Microclimatic, and Open-Quantum Assessment of Semi-Transparent Organic Agrivoltaics for Sub-Saharan Perennial Crops''} for consideration in \textit{Renewable Energy}.

Wavelength-selective semi-transparent organic photovoltaics (ST-OPVs) are widely proposed for tropical agrivoltaics on the premise that transmitted spectra enhance crop productivity relative to opaque crystalline silicon (c-Si) while modulating photosynthetic quantum coherence. Yet existing evaluations rarely link calibrated transfer-matrix optics, hourly tropical microclimate and stomatal--soil-water physiology, cooperative processing dispatch, replacement-aware life-cycle accounting, and global uncertainty propagation within one self-consistent counterfactual framework. Here, we assess three flagship perennial production chains across western Cameroon's wet--dry highland belt---Arabica coffee (Foumbot, 1,044~m), Protected Geographical Indication (PGI) Penja white pepper (Njombe-Penja, 273~m), and highland tea (Djuttitsa, 1,400~m)---over a full 8,760~h (2023 ERA5) meteorological cycle.

Exact four-lever Shapley attribution reveals a pronounced seasonal regime inversion: dry-season leaf cooling ($-1.5$ to $-3.5\,^\circ\mathrm{C}$) boosts dry-season assimilation across all species ($+8.8\%$ to $+30.2\%$ rainfed; $+49.7\%$ irrigated), whereas monsoon cloudiness drives net annual yields negative in coffee ($-7.08\%$) and tea ($-3.35\%$) under $40.6\%$ shade, leaving Penja white pepper as the sole crop sustaining net rainfed gains ($+7.43\%$ to $+8.16\%$, land equivalent ratio $1.48$--$1.55$). Crucially, at matched $\sim 40\%$ shade, ST-OPV and elevated c-Si crop yields differ by at most $1.05$ percentage points, demonstrating that shade fraction and seasonal timing govern agronomy. Furthermore, multi-Hamiltonian continuous-wave steady-state bounds (validated against MesoHOPS~1.8.0 and exact HEOM) constrain excitonic efficiency shifts below $3.37\times 10^{-5}$ ($2.10\times 10^{-3}$ HEOM envelope), resulting in negligible Sobol sensitivity ($S_{Ti} \le 2.60\times 10^{-7}$).

Across $N=10,000$ Saltelli quasi-Monte Carlo realizations, dual variance decomposition shows that grid export tariffs and cost of capital govern standalone ST-OPV bankability, whereas hardware costs ($C_{\mathrm{mod}}^{\mathrm{OPV}}, C_{\mathrm{BOS}}^{\mathrm{Si}}$) and encapsulation lifetime ($T_{80}$) govern competition against elevated c-Si, which achieves higher 20-year NPV in $77.0\%$--$89.0\%$ of draws ($75.3\%$--$87.7\%$ at zero carbon). Life-cycle carbon parity requires $T_{80,\mathrm{LCA}}^* \ge 7.40$--$8.12$~yr (slashing intensity to $7.8$--$12.5~\mathrm{gCO_2e/kWh}$ at 25~yr), whereas economic NPV parity emerges only in Penja pepper at $\$15/\mathrm{m}^2$ with $T_{80,\mathrm{NPV}}^* \ge 16.33$~yr; overnight tea withering storage lifts solar fraction to $99.9\%$ and sustains positive off-grid NPV ($+\$319,794/\mathrm{ha}$).

This manuscript is original, has not been published previously, and is not under consideration elsewhere. All authors approve this submission and declare no competing financial interests.

\vspace{0.6em}
\noindent
Sincerely,

\vspace{0.5em}
\noindent
Serge Guy Nana Engo (Corresponding author, on behalf of all co-authors)\\
Department of Physics, Faculty of Science, University of Yaound\'e~I, P.O. Box 812, Yaound\'e, Cameroon\\
E-mail: serge.nana-engo@facsciences-uy1.cm
\end{document}
